\section{Training}
\label{sec:training}

\subsection{SFT Recipe Overview}

AmphionASR is initialised from the public
Qwen3-ASR-1.7B weights~\cite{qwen2025qwen3asr} and produced by a
three-stage supervised fine-tuning (SFT) recipe with all-linear LoRA
over the eight training sets of Table~\ref{tab:training_sets}.
LoRA uses rank $64$ and $\alpha=128$, applied to every linear
projection of each trainable module. All stages use
ms-swift~\cite{msswift} with DeepSpeed ZeRO-1; each stage runs for
one epoch at learning rate $1\!\times\!10^{-6}$ with cosine decay and
a $5$\,\% warmup ratio, effective batch size $384$ on six GPUs, and
$1$\,\% of examples held out for validation.

\paragraph{Stage~1 -- degradation and whisper; encoder and LLM jointly tuned.}
The audio encoder and the LLM are both trainable. Training uses the
Voice-in-the-Wild-2M degradation set~\cite{xie2026mega} and the
WhispEar whispered-speech set~\cite{fang2025whisperear}.
Both domains contain acoustic conditions that differ from clean
speech.
We jointly update the audio encoder and the LLM to allow both modules
to adapt before prompt-conditioned training is introduced.

\paragraph{Stage~2 -- hotword and target-speaker; encoder frozen.}
The audio encoder is frozen and only the LLM is updated. Training uses
the hotword-conditioned mixture from five public corpora
(Table~\ref{tab:data_hw_train}) together with three synthesised
target-speaker partitions (Table~\ref{tab:ts_training_sets}).
This stage trains the LLM on hotword and target-speaker conditions
while retaining the Stage~1 encoder.

\paragraph{Stage~3 -- replay against catastrophic forgetting; encoder frozen.}
The audio encoder remains frozen and only the LLM is updated. The
training pool combines the Stage~1 and Stage~2 training sets with
two plain-transcription corpora replayed to mitigate catastrophic
forgetting: AISHELL-2~\cite{du2018aishell2} and
LibriSpeech~\cite{panayotov2015librispeech}.
The replay mixture includes both conditioned and unconditioned
examples to reduce the risk of forgetting plain transcription while
continuing training on the Stage~2 tasks.

\subsection{Hotword Retriever}
\label{subsec:hotword-retriever}

The hotword retriever is trained separately from the three-stage SFT
recipe, on top of the merged Stage~2 checkpoint, optimising only the
two small MLP adapters of Section~\ref{subsec:hotword-rag}; the audio
encoder and the LLM token-embedding table remain frozen. Training is
bilingual with per-language hotword pools: we collect the hotwords
annotated in the CommonVoice Mandarin and English hotword splits,
deduplicate them, and build a Mandarin pool and an English pool that
serve as the candidate vocabulary for each language.

We build the retrieval objective upon
GLCLAP~\cite{kong2025glclap}, adopting its global--local bidirectional
contrastive loss: a global term aligning pooled audio with the
transcript and a local term aligning each hotword with per-frame audio
via a max-over-time similarity, each summed over the
speech$\to$text and text$\to$speech directions. Our only deviation is
the positive/negative selection of the local branch. The positive for
each utterance is its ground-truth hotword; instead of GLCLAP's
in-batch negatives, we sample $N{=}4095$ negatives per utterance from the
same-language hotword pool, excluding the batch's own hotwords and
shared across the batch so that the negative embeddings are encoded
only once per step. The text$\to$speech direction keeps in-batch audios
as candidates. The temperature is learnable.
Training runs for $10$ epochs at learning rate $3\!\times\!10^{-4}$
with cosine decay, a $5$\,\% warmup ratio, a learnable temperature
initialised at $0.07$, gradient clipping at $1.0$, and effective batch
size $384$ on six GPUs under fp16. The checkpoint with the best
validation Recall@$1$ on the CommonVoice hotword splits is shipped as
the retrieval module.

\subsection{Group Relative Policy Optimisation}
\label{subsec:grpo}

After the three-stage SFT recipe, we apply Group Relative Policy
Optimisation (GRPO) to the language backbone while keeping the audio
encoder frozen.
The sequence-level reward combines transcription accuracy with
agreement on whether each candidate hotword appears in the output.

\paragraph{Training data.}
GRPO uses a subset of the hotword-conditioned and plain-transcription
SFT data, mixed at a $4{:}1$ ratio. Sampling preserves the task and
language proportions of the corresponding SFT data.

Data preprocessing and prompt construction otherwise follow the SFT
recipe; the GRPO-specific optimisation settings are reported below.

\paragraph{Rewards.}
The accuracy reward is the negative CER clipped to $[0, 1]$,
\begin{equation}
  R_\mathrm{asr}(\hat{y}, y) = \max(0,\, 1 - \cer(\hat{y}, y)),
  \label{eq:reward-asr}
\end{equation}
with empty-reference boundary cases scored as full credit.
The hotword match-accuracy reward averages per-candidate binary
agreement between predicted and reference hotword presence,
\begin{equation}
  R_\mathrm{hw}(\hat{y}, y; C) = \frac{1}{|C|}\sum_{c \in C}
    \mathbf{1}\!\left\{\mathbf{1}[c \in \mathrm{Pred}] = \mathbf{1}[c \in \mathrm{Ref}]\right\}.
  \label{eq:reward-hw}
\end{equation}
The total reward is the weighted sum
\begin{equation}
  R = w_\mathrm{asr}\, R_\mathrm{asr} + w_\mathrm{hw}\, R_\mathrm{hw},
  \qquad (w_\mathrm{asr}, w_\mathrm{hw}) = (1.0, 0.3).
  \label{eq:reward-total}
\end{equation}
The GRPO run uses group size $G=16$, sampling temperature $0.6$, PPO
clip $\varepsilon=0.2$, KL coefficient $\beta=1\!\times\!10^{-3}$,
peak learning rate $1\!\times\!10^{-6}$ for one epoch, and a
per-device prompt batch of $8$ with gradient accumulation of $2$.

The evaluation reports the resulting model as a complete training
recipe. Without an SFT-only control, it does not isolate the
contribution of GRPO from the preceding supervised stages.
